% ============================================================== 2. Related work
\section{Related work}
\label{sec:related}

\subsection{Federated learning for medical imaging and dermatology}

FedAvg \cite{mcmahan2017fedavg} trains a global model by averaging locally updated parameters. Because
clients hold non-identically distributed data, many variants correct client drift, for example FedProx
with a proximal term \cite{li2020fedprox} and SCAFFOLD with control variates
\cite{karimireddy2020scaffold}, or personalize parts of the model, for example local batch-normalization
layers in FedBN \cite{li2021fedbn}, local classification layers in FedPer \cite{arivazhagan2019fedper},
or local fine-tuning of the global model; surveys review these
families \cite{kairouz2021advances,tan2023pfl}. In medical imaging, FL has been studied for brain tumor
segmentation, chest imaging and outcome prediction across hospitals
\cite{sheller2020federated,dayan2021exam}, and its promises and pitfalls have been reviewed
\cite{rieke2020future}. FLamby introduced Fed-ISIC2019 as a cross-silo dermatology
benchmark with clients defined by acquisition source \cite{flamby2022}, and Tavakoli et al.\ used it to
compare personalized methods on heterogeneous clinical data \cite{tavakoli2024fendafl}. Both evaluations
use FLamby's fixed image-level split. Dermatology foundation models such as PanDerm, pretrained on
millions of skin images, now offer strong frozen features \cite{yan2025panderm}; how leakage affects
comparisons on such features has not been examined.

\subsection{Data leakage and evaluation pitfalls in medical imaging}

Leakage, in which information about the test data reaches the training process, is a widespread cause
of irreproducible results in machine-learning-based science \cite{kapoor2023leakage}. In medical imaging
it often arises when images from the same subject appear in training and test sets: splitting
histopathology slides at the tile level \cite{bussola2021tiles} and brain MRI volumes at the slice level
\cite{yagis2021leakage} markedly overestimated classification performance compared with subject-level
splits. Related failure modes include confounding by acquisition site \cite{zech2018confounding} and shortcut
learning \cite{degrave2021shortcuts}, both of which
can make held-out performance unrepresentative \cite{roberts2021pitfalls,varoquaux2022failures}. In
dermatology, Cassidy et al.\ found duplicates across the yearly ISIC releases and proposed curated
training and test sets \cite{cassidy2022isic}; Gr\"oger et al.\ removed near-duplicates and other
problems from six dermatology evaluation benchmarks \cite{groger2023reliable}; and Abhishek et al.\ showed
that duplicates leaked across the splits of DermaMNIST and Fitzpatrick17k and released corrected
partitions \cite{abhishek2025dermamnist}. These studies establish that dermatology benchmarks can contain
leakage; they concern centralized datasets and accuracy. What we add is the per-client distribution of
lesion-level leakage in a federated benchmark, a randomized estimate of its effect, and its consequences
for method rankings, calibration and conformal coverage.

\subsection{Calibration of medical image classifiers}

Modern neural networks are often over-confident \cite{guo2017calibration}. Calibration is typically
measured with the expected calibration error \cite{naeini2015ece}, class-wise variants
\cite{kull2019beyond}, the negative log-likelihood and the Brier score \cite{brier1950}, and improved post
hoc by temperature scaling \cite{guo2017calibration} or by approximate Bayesian and ensemble methods
\cite{gal2016dropout}. Calibration degrades under dataset shift
\cite{ovadia2019trust} and depends on architecture and training \cite{minderer2021revisiting}, and it has
been studied in medical imaging, for example for segmentation networks \cite{mehrtash2020confidence}. The question we
address is narrower: how calibration varies across the clients of Fed-ISIC2019 and whether a test set with
lesion-level leakage changes calibration estimates and the effect of temperature scaling.

\subsection{Conformal prediction in federated settings}

Split conformal prediction turns any classifier into a set-valued predictor with marginal coverage of at
least $1-\alpha$ when calibration and test points are exchangeable
\cite{vovk2005algorithmic,angelopoulos2023gentle}; common scores for classification
are LAC \cite{sadinle2019lac} and APS \cite{romano2020aps}. Coverage can fail under
covariate or label shift \cite{tibshirani2019covariate,podkopaev2021labelshift,barber2023beyond}, can
differ between subgroups \cite{lu2022fair}, and needs adapted arguments when observations are grouped,
as are repeated images of one lesion \cite{dunn2022hierarchical,lee2023hierarchical}. Several coverage
targets must be distinguished: marginal coverage on the mixture of clients, coverage for each client,
coverage for each class, and the probability over calibration samples that coverage falls short.
Federated conformal methods compute thresholds without pooling scores: Lu et al.\ derived a federated
quantile with a guarantee on the mixture of clients under partial exchangeability
\cite{lu2023fcp}, Humbert et al.\ proposed a one-shot quantile-of-quantiles scheme
\cite{humbert2023oneshot}, Plassier et al.\ addressed label shift between clients
\cite{plassier2023labelshift}, Min et al.\ personalized conformal sets through localization
\cite{min2025pfcp}, and Wen et al.\ targeted group-conditional coverage \cite{wen2026gcfcp}. A method can
meet its mixture target while under-covering a particular client; it is not disproved by failing a per-client objective it was not designed for, but it is insufficient for site-level use. Per-client thresholds avoid that bias
but are noisy for small clients. Shrinking noisy estimates toward a common mean is a classical idea,
from normal-means empirical Bayes \cite{efron1975stein} to random-effects meta-analysis
\cite{dersimonian1986meta}. EB-ShrinkFCP applies it to per-client conformal quantiles with only two
numbers per client; its methodological novelty is limited to this application, and we evaluate it
empirically on real clients, on held-out clients and against client-wise, global, FCP and
label-shift-weighted thresholds.
